\section{Conclusion}
\label{sec:conclusion}

LLM repair agents run tests indiscriminately and most of that execution does
not pay~\cite{lin2026run,chen2026rethinking}. Classical repair cuts that work
either by relating candidate patches to one
another~\cite{weimer2013leveraging,mechtaev2018testequivalence}, which an
unenumerable candidate space rules out, or by ordering the evidence so a doomed
candidate dies on the first
case~\cite{qi2013efficient,venugopal2020modification}.
This paper carries the second mechanism into an LLM loop and changes what it
orders from: counterexamples retained across candidates, where an input already
in memory \emph{proves} that a candidate reproducing it will fail. Memory cuts
oracle calls \fold{5.2} at unchanged repair rate and, counting every guard
replay as the execution it is, total program executions \fold{2.2}. That there
is anything to intercept is itself a finding: one proposal in four is
byte-identical to an earlier one, a rate that is zero by construction in the
synthesis loop the design came from.

Two results change the recommendation. Steering does nothing at this scale, at
$+128\%$ prompt tokens, and makes the full agent the slowest of five arms; and
the $\loc$ \emph{ordering} is indistinguishable from a flat scan on memories
this small---the typed guard's $39\%$ cheaper seconds come from deduplication,
not the index. What we recommend is therefore not the
architecture we set out to validate but the guard alone. It does not survive a per-modification-point prioritizer on
executions---that prioritizer is cheaper on $80$ of $99$ tasks, by a median
$4.6\%$---and it is the only policy of five that removes oracle invocations, on
every task (\cref{sec:results-policies}). Whether that holds at another
model scale, or on the expensive-oracle faults we filtered out, are the two
questions we would want answered next.
